\subsubsection*{\texttt{30-retrieval-bm25}}
\label{app:samples:30_retrieval_bm25}

\paragraph{Q3: How long does the Haste spell last?}
\noindent \textbf{Corr.:}~5.00/5\quad \textbf{Faith.:}~1.00\quad \textbf{Relev.:}~1.00\quad \textbf{Ctx.Rel.:}~1.00\\

\noindent\textbf{Antwort:} The Haste spell lasts for a duration of concentration, up to 1 minute.

\noindent\textbf{Abgerufene Chunks:}
\begin{itemize}[noitemsep, topsep=2pt]
  \item[\textbf{\#1}] \texttt{[0.0333]} \textit{Spells/Haste} -- \#\#\# Haste  *3rd-level transmutation*  **Casting Time:** 1 action  **Range:** 30 \ldots
  \item[\textbf{\#2}] \texttt{[0.0320]} \textit{Spells/Modify Memory} -- \#\#\# Modify Memory  *5th-level enchantment*  **Casting Time:** 1 action  **Range:\ldots
  \item[\textbf{\#3}] \texttt{[0.0164]} \textit{Spells/# Spellcasting} -- \#\# Duration  A spell's duration is the length of time the spell persists. A dura\ldots
  \item[\textbf{\#4}] \texttt{[0.0164]} \textit{Gamemastering/Traps} -- \#\# Traps in Play  When adventurers come across a trap, you need to know how the \ldots
  \item[\textbf{\#5}] \texttt{[0.0161]} \textit{Gamemastering/Madness} -- \#\# Madness Effects  Madness can be short-term, long-term, or indefinite. Most re\ldots
\end{itemize}
\vspace{4pt}
\hrule
\vspace{6pt}

\paragraph{Q5: Can a grappled creature cast spells?}
\noindent \textbf{Corr.:}~1.00/5\quad \textbf{Faith.:}~1.00\quad \textbf{Relev.:}~0.00\quad \textbf{Ctx.Rel.:}~0.25\\

\noindent\textbf{Antwort:} The context does not provide information on whether a grappled creature can cast spells or not. The first context only mentions that a grappled creature's speed becomes 0 and it can't benefit from any bonus to its speed. The second context is about spellcasting for monsters, but it doesn't mention a\ldots

\noindent\textbf{Abgerufene Chunks:}
\begin{itemize}[noitemsep, topsep=2pt]
  \item[\textbf{\#1}] \texttt{[0.0323]} \textit{Gamemastering/Conditions} -- \#\# Grappled  - A grappled creature's speed becomes 0, and it can't benefit from \ldots
  \item[\textbf{\#2}] \texttt{[0.0323]} \textit{Gameplay/Combat} -- \#\#\# Grappling  When you want to grab a creature or wrestle with it, you can use \ldots
  \item[\textbf{\#3}] \texttt{[0.0315]} \textit{Spells/Arcane Hand} -- ***Grasping Hand***. The hand attempts to grapple a Huge or smaller creature wit\ldots
  \item[\textbf{\#4}] \texttt{[0.0167]} \textit{Gameplay/Combat} -- \#\# Cast a Spell  Spellcasters such as wizards and clerics, as well as many monst\ldots
  \item[\textbf{\#5}] \texttt{[0.0164]} \textit{Monsters/# Monster Statistics} -- \#\#\# Innate Spellcasting  A monster with the innate ability to cast spells has th\ldots
\end{itemize}
\vspace{4pt}
\hrule
\vspace{6pt}

\paragraph{Q12: If a creature has exhaustion level 3 and is also restrained, what is its movement speed?}
\noindent \textbf{Corr.:}~1.00/5\quad \textbf{Faith.:}~0.00\quad \textbf{Relev.:}~0.00\quad \textbf{Ctx.Rel.:}~1.00\\

\noindent\textbf{Antwort:} I cannot find the rules in the provided context regarding a creature's movement speed when it has exhaustion level 3 and is also restrained. The context only mentions that exhaustion affects attack rolls and Wisdom (Perception) checks that rely on sight. The Haste spell description does not mention \ldots

\noindent\textbf{Abgerufene Chunks:}
\begin{itemize}[noitemsep, topsep=2pt]
  \item[\textbf{\#1}] \texttt{[0.0328]} \textit{Gamemastering/Conditions} -- \#\# Exhaustion  Some special abilities and environmental hazards, such as starvat\ldots
  \item[\textbf{\#2}] \texttt{[0.0325]} \textit{Gamemastering/Conditions} -- \#\# Restrained  - A restrained creature's speed becomes 0, and it can't benefit f\ldots
  \item[\textbf{\#3}] \texttt{[0.0323]} \textit{Spells/Freedom of Movement} -- \#\#\# Freedom of Movement  *4th-level abjuration*  **Casting Time:** 1 action  **R\ldots
  \item[\textbf{\#4}] \texttt{[0.0164]} \textit{Gameplay/Adventuring} -- \#\# Suffocating  A creature can hold its breath for a number of minutes equal to \ldots
  \item[\textbf{\#5}] \texttt{[0.0161]} \textit{Gameplay/Combat} -- \#\# Flying Movement  Flying creatures enjoy many benefits of mobility, but they m\ldots
\end{itemize}
\vspace{4pt}
\hrule
\vspace{6pt}

\paragraph{Q16: Can you use Action Surge to cast two leveled spells in one turn?}
\noindent \textbf{Corr.:}~2.00/5\quad \textbf{Faith.:}~1.00\quad \textbf{Relev.:}~1.00\quad \textbf{Ctx.Rel.:}~0.50\\

\noindent\textbf{Antwort:} I cannot find this in the provided rules about using Action Surge to cast two leveled spells in one turn for a Fighter or any other class. Action Surge only allows you to take one additional action on top of your regular action and a possible bonus action. Casting a spell is considered an action, so\ldots

\noindent\textbf{Abgerufene Chunks:}
\begin{itemize}[noitemsep, topsep=2pt]
  \item[\textbf{\#1}] \texttt{[0.0333]} \textit{Classes/Fighter} -- \#\#\# Action Surge  Starting at 2nd level, you can push yourself beyond your norma\ldots
  \item[\textbf{\#2}] \texttt{[0.0315]} \textit{Spells/# Spellcasting} -- \#\#\# Bonus Action  A spell cast with a bonus action is especially swift. You must\ldots
  \item[\textbf{\#3}] \texttt{[0.0303]} \textit{Classes/Sorcerer} -- \#\#\#\# Flexible Casting  You can use your sorcery points to gain additional spell \ldots
  \item[\textbf{\#4}] \texttt{[0.0164]} \textit{Classes/Fighter} -- **Table- The Fighter**  | Level | Proficiency Bonus | Features                  \ldots
  \item[\textbf{\#5}] \texttt{[0.0161]} \textit{Classes/Sorcerer} -- \#\#\#\# Quickened Spell  When you cast a spell that has a casting time of 1 action,\ldots
\end{itemize}
\vspace{4pt}
\hrule
\vspace{6pt}
